\documentclass[10pt,letterpaper]{article}

\usepackage[utf8]{inputenc}
\usepackage[T1]{fontenc}
\usepackage{graphicx}
\usepackage{amsmath,amssymb}
\usepackage{hyperref}
\usepackage{url}
\usepackage{cite}
\usepackage{booktabs}
\usepackage[margin=1in]{geometry}
\usepackage{microtype}
\usepackage{textcomp}
\usepackage{authblk}

\title{Agent-based modeling reveals nonlinear, synergistic cytokine interactions that shape skeletal muscle regeneration}

\author[1]{Author Placeholder}
\affil[1]{Department Placeholder, Institution Placeholder}

\date{}

\begin{document}
\maketitle

\begin{abstract}
Skeletal muscle regeneration after acute injury is an emergent process in which cytokines, immune cells, satellite stem cells (SSCs), fibroblasts, and the microvasculature interact across space and time. Combinatorial manipulation of cytokine dynamics is experimentally infeasible because of the number of possible combinations and the cost of longitudinal cell-resolved readouts. To bridge this gap, we developed a spatial agent-based model (ABM) of murine skeletal muscle regeneration implemented in the Cellular-Potts framework. The model represents a two-dimensional fascicle cross-section of approximately fifty muscle fibers, together with neutrophils, monocytes, resident and polarized macrophages, SSCs, fibroblasts, myofibroblasts, myoblasts, myocytes, capillaries, and lymphatic vessels, and it tracks seven diffusible cytokines whose diffusivity varies with local collagen density. We curated over one hundred literature-derived rules and calibrated fifty-two unknown parameters using the CaliPro parameter density estimation protocol against temporal datasets for cross-sectional area (CSA), SSC, and fibroblast counts. Held-out outputs for macrophage, neutrophil, and capillary counts tracked experimental trends without retuning, and the model qualitatively reproduced eight published perturbation experiments spanning VEGF-A injections, cytokine knockouts, cell depletions, and hindered angiogenesis. Latin hypercube sampling combined with partial rank correlation coefficient (PRCC) analysis identified cytokine parameters whose alteration improved regeneration; the optimal combination of cytokine decay and diffusion changes produced a 13\% improvement in CSA recovery at 28 days compared with unaltered regeneration, exceeding the benefit of any single-parameter intervention. These results illustrate how a spatially resolved ABM can expose nonlinear thresholds and synergies that guide the design of combinatorial cytokine therapies for muscle injury.
\end{abstract}

%========================================================================
\section{Introduction}
\label{sec:introduction}

Acute skeletal muscle injuries are among the most common orthopedic complaints, accounting for a substantial fraction of all musculoskeletal injuries seen in athletic and clinical populations \cite{maffulli2014muscle}. Standard of care is largely supportive---rest, ice, compression, elevation, and anti-inflammatory medication---and these interventions rest on a thin mechanistic foundation \cite{deyhle2018role}. Outcomes are inconsistent: more than 30\% of reported orthopedic injuries involve muscle, and incomplete functional recovery, scar formation, and high recurrence rates remain common \cite{maffulli2014muscle}. The underlying difficulty is biological, not simply clinical. Healthy muscle regeneration is an emergent, multistage cascade in which coordinated interaction among cells and their microenvironment, rather than any single cell type, determines whether recovery is complete, delayed, or fibrotic \cite{forcina2020mechanisms,chazaud2020inflammation}.

The regeneration cascade unfolds over weeks and involves a time-ordered sequence of events. Within hours of injury, neutrophils and monocytes are recruited through perfused capillaries to clear necrotic tissue, while resident macrophages and newly arrived monocytes differentiate into pro-inflammatory (M1) macrophages that phagocytose debris and secrete signals that activate fibroblasts and SSCs \cite{tidball2017regulation,dumont2015satellite}. As inflammation resolves, M1 macrophages transition to an anti-inflammatory (M2) phenotype, and activated SSCs proliferate, differentiate into myoblasts, and fuse with one another or with damaged fibers to rebuild the cross-sectional area \cite{arnold2007inflammatory,bentzinger2013cellular}. Microvascular remodeling proceeds in parallel: damaged capillaries become non-perfused, and new sprouts form under angiogenic stimulation driven by vascular endothelial growth factor A (VEGF-A) and matrix metalloproteinase-9 (MMP-9) \cite{carmeliet2003angiogenesis}. Throughout this choreography, cytokines including hepatocyte growth factor (HGF), transforming growth factor beta (TGF-$\beta$), tumor necrosis factor alpha (TNF-$\alpha$), MMP-9, VEGF-A, monocyte chemoattractant protein-1 (MCP-1), and interleukin 10 (IL-10) act as the signaling backbone that coordinates cell decisions \cite{tatsumi1998hgf,lu2011macrophages}.

Although the individual roles of these cytokines have been characterized in isolation, systematically probing the combinatorial space of cytokine manipulations is experimentally intractable. Realistic variations of decay and diffusion parameters across seven cytokines generate more than a million distinct parameter combinations, and each experimental arm requires a reproducible injury model, timepoint-resolved tissue sampling, and dedicated immunoassays \cite{wehling2009interleukin}. Pleiotropic actions, soluble receptors, and cross-regulation further confound direct interpretation. These barriers motivate an \emph{in silico} approach capable of exploring the combinatorial landscape, generating mechanistic hypotheses, and prioritizing experiments. Prior computational work has laid important groundwork in this direction \cite{virgilio2015multiscale,martin2017agent}, but most published models of muscle have relied on simplified, non-spatial cytokine representations and static vasculatures, limiting their ability to capture transport-based interventions or angiogenic dynamics.

In this work, we develop a spatial agent-based model (ABM) of murine skeletal muscle regeneration implemented in the Cellular-Potts framework \cite{graner1992simulation,swat2012compucell3d}. The model represents muscle fibers, satellite stem cells, fibroblasts, neutrophils, macrophages, microvessels, and lymphatic vessels, together with seven diffusing cytokines whose dynamics depend on local collagen density. We calibrate unknown parameters by parameter density estimation using the CaliPro protocol \cite{joslyn2021calipro}, validate held-out outputs against published experimental data, and reproduce eight perturbation experiments drawn from the literature. Finally, we use Latin hypercube sampling combined with partial rank correlation coefficient analysis \cite{marino2008sensitivity} to identify cytokine decay and diffusion parameters whose alteration enhances CSA recovery. The goals of the study are to (i) build a spatial ABM that jointly captures cytokine diffusion and microvascular adaptation, (ii) calibrate and validate against multiple independent datasets, and (iii) use the calibrated model to generate testable hypotheses about cytokine interventions that could improve regeneration outcomes.

%========================================================================
\section{Related Work}
\label{sec:related}

Computational modeling of skeletal muscle regeneration has advanced rapidly over the last decade. Early non-spatial ordinary differential equation (ODE) models captured bulk population dynamics but could not resolve the spatial gradients of cytokines that shape local cell decisions \cite{hillen2009user}. Foundational agent-based models of muscle demonstrated that cell-level rules can reproduce tissue-scale behaviors in healthy muscle and in Duchenne muscular dystrophy, and that \emph{in silico} perturbations can be used to test therapeutic hypotheses \cite{virgilio2015multiscale,martin2017agent}. However, these implementations relied on simplified, non-spatial cytokine representations. As a result, prior ABMs could not capture interventions that target cytokine transport (for example, injection of recombinant TGF-$\beta$ or use of plasmin activators to release ECM-bound VEGF-A), nor the feedback between cytokine distributions and local ECM remodeling \cite{urciuolo2013collagen}.

A second limitation of earlier muscle ABMs has been the absence of an adaptive microvasculature. Capillary perfusion state, angiogenic sprouting, and the balance between VEGF-A and MMP-9 are central to the recovery trajectory \cite{carmeliet2003angiogenesis}, yet most published ABMs of muscle treat the vasculature as a static boundary condition rather than as adaptive agents that respond to injury, hypoxia, and cytokine gradients. The lymphatic system and its role in cytokine clearance and cell migration have been similarly neglected. Recent ABMs have begun to address cerebral palsy and eccentric contraction injuries, but angiogenic remodeling during acute regeneration remains underrepresented in the ABM literature.

The Cellular-Potts framework, originally introduced for biological cell sorting \cite{graner1992simulation} and implemented in the CompuCell3D platform \cite{swat2012compucell3d}, is well suited to bridge these gaps because it represents cells as extended lattice domains whose morphology and motility arise from a free-energy minimization and whose surrounding fields (chemokines, ECM) are solved on the same lattice. Combining this representation with CaliPro-based parameter density estimation \cite{joslyn2021calipro} and LHS/PRCC global sensitivity analysis \cite{marino2008sensitivity} provides a route from literature-derived rules to calibrated, validated, and interpretable predictions. Our model builds on this foundation by (a) introducing explicit cytokine-specific diffusion that responds to local collagen density, (b) adding dynamic capillary and lymphatic agents capable of perfusion loss, angiogenic sprouting, and cytokine uptake, and (c) using sensitivity analysis to design \emph{in silico} perturbations that would be infeasible experimentally.

%========================================================================
\section{Methods}
\label{sec:methods}

\subsection{Model architecture and agents}

The ABM represents a two-dimensional cross-section of murine skeletal muscle containing approximately fifty fibers. The domain geometry is extracted from a laminin-stained histology image, masked to separate fibers from ECM, and converted into a CompuCell3D initialization script that labels pixels as muscle fibers, ECM, capillaries, or lymphatic vessels. Simulations run for 28 days of biological time; each Monte Carlo step (mcs) corresponds to 15 minutes and incorporates 45 Cellular-Potts evaluations per step so that emergent cell speeds match literature-derived values. The full agent inventory comprises muscle fibers, necrotic fibers, ECM, capillaries, lymphatic vessels, quiescent and activated fibroblasts, myofibroblasts, quiescent and activated SSCs, myoblasts, myocytes, immature myotubes, neutrophils, monocytes, resident macrophages, M1 macrophages, and M2 macrophages. Seven diffusing factors are tracked: HGF, MCP-1, MMP-9, TGF-$\beta$, TNF-$\alpha$, VEGF-A, and IL-10. A schematic of the full interaction structure is shown in Figure~\ref{fig:overview}.

\begin{figure}[t]
  \centering
  \includegraphics[width=0.95\linewidth]{figures/fig_overview.pdf}
  \caption{Overview of the agent-based model. Cells, cytokines, and the extracellular matrix interact through recruitment, chemotaxis, phagocytosis, secretion, and binding, with cytokine diffusion recalculated as local collagen density changes.}
  \label{fig:overview}
\end{figure}

\subsection{Agent rules and cytokine dynamics}

Agent behaviors are governed by literature-derived, probability-weighted rules. Neutrophils are recruited through perfused capillaries, chemotax up HGF gradients, phagocytose necrotic tissue, secrete MMP-9, MCP-1, and TNF-$\alpha$ during phagocytosis, and apoptose after two phagocytic events or a maximum residence time. Monocytes transition to M1 macrophages based on cumulative TNF-$\alpha$ exposure or a sampled transition time, and M1 macrophages transition to M2 macrophages when they exceed an IL-10 concentration threshold or accumulate sufficient phagocytic events. SSCs follow a more elaborate cycle with 33 parameters across 17 rules, capturing quiescence, activation by HGF, symmetric and asymmetric division, differentiation into myoblasts and myocytes, and apoptosis driven by TGF-$\beta$ in the absence of VEGF-A or protective macrophage neighbors. Fibroblasts are activated by TGF-$\beta$, chemotax toward regions of low-density collagen using a spring-like potential, secrete collagen at ECM sites, and differentiate into myofibroblasts on prolonged TGF-$\beta$ exposure.

For each cytokine-dependent behavior, binding and uptake are coupled to removal of the cytokine from the field, simulating ligand-receptor binding and endocytosis. Because direct measurements of tissue-context cytokine diffusivity are limited, we estimate effective diffusivity using the collagen-volume-fraction relationship proposed by prior work \cite{urciuolo2013collagen}. The collagen fraction is updated throughout the simulation as fibroblasts and myofibroblasts deposit collagen and as MMP-9 released by phagocytes degrades it; cytokine diffusivity is recalculated accordingly at each mcs.

\subsection{Calibration with CaliPro}

Of the parameters required to close the model, many are drawn from published measurements. Fifty-two parameters lack direct physiological measurement and must be inferred from data. To do this, we applied CaliPro \cite{joslyn2021calipro}, a model-agnostic calibration protocol that uses parameter density estimation and alternative density subtraction (ADS) to iteratively narrow parameter ranges. We generated 600 Latin hypercube samples across the candidate parameter space, ran each in triplicate, and classified runs as passing or failing relative to predefined criteria. Density functions of passing and failing runs were computed in each dimension, and new bounds were set where the passing density dominates. The calibration targets were time-varying CSA, SSC counts, and fibroblast counts, each normalized by the experimental peak or a designated reference day (day 5 for SSCs and fibroblasts, day 1 for neutrophils, day 3 for macrophages). Passing criteria were that the model 95\% confidence interval fell within 1 SD of experimental CSA and within 2.5 SD of SSC and fibroblast counts. Eight iterations of ADS narrowed the space until further tightening eliminated viable parameters, at which point we reverted to the previous iteration and fixed bounds so that all three replicates of the selected parameter set lay within the strictest criteria. The final parameter set was run 100 times to verify stochastic reproducibility.

\subsection{Validation and perturbation studies}

Validation used held-out cell counts that were not part of the calibration targets. M1, M2, and total macrophage counts, neutrophil counts, and capillary counts were compared with published experimental data without further tuning. In addition, eight perturbation conditions were run and compared against corresponding experimental findings: four levels of VEGF-A injection (low, medium, high, extra high at relative concentrations of 500, 750, $10^3$, and $2 \times 10^3$), hindered angiogenesis in which damaged capillaries could not reperfuse, IL-10 knockout, MCP-1 knockout, TNF-$\alpha$ knockout, and macrophage depletion. Each perturbation was run for 100 replicates, and outputs were compared with control simulations using a two-sample $t$-test with a significance level of 0.05.

\subsection{Sensitivity analysis and \emph{in silico} experiments}

Global sensitivity analysis used Latin hypercube sampling combined with partial rank correlation coefficients \cite{marino2008sensitivity}. We sampled cytokine decay and diffusion coefficients for all seven cytokines across a range from 0.1 to 10 times their calibrated values while holding other parameters fixed, generating 300 parameter sets run in triplicate. PRCCs were computed every 10 hours for CSA and for cell counts of SSCs, fibroblasts, non-perfused capillaries, myoblasts, myocytes, neutrophils, M1, and M2 macrophages. We applied an $\alpha = 0.05$ threshold with a Bonferroni correction for the number of tests. The PRCC profile was used to select cytokine parameter perturbations that individually promoted beneficial regeneration metrics, and these individual perturbations were then combined to assess whether their effects were additive or synergistic.

%========================================================================
\section{Results}
\label{sec:results}

\subsection{Calibration and held-out validation}

After eight iterations of density-based calibration, the final parameter set produced simulated CSA, SSC, and fibroblast trajectories that tracked the experimental means and fell within the specified SD envelopes at nearly every timepoint. The only target not fully captured by the model 95\% confidence interval was the SSC count at day 3, which fell slightly outside the experimental SD. For the remaining calibration targets, the model 95\% confidence interval was within the experimental SD at every timepoint. These results indicate that the parameter density estimation procedure converged to a region of parameter space that jointly satisfies multiple temporal datasets despite their differing sources and experimental conditions. Figure~\ref{fig:calibration} summarizes this outcome across calibration and held-out targets. Importantly, the held-out outputs---total macrophage, M1, M2, neutrophil, and capillary counts---were consistent with their respective experimental trends without any retuning, providing independent evidence that the calibrated parameter set encodes realistic cytokine-cell dynamics rather than overfitting the calibration targets.

\begin{figure}[t]
  \centering
  \includegraphics[width=0.90\linewidth]{figures/fig_calibration.pdf}
  \caption{Calibration and held-out validation outcomes. Green bars indicate targets whose model 95\% confidence interval falls within the experimental standard deviation; the single exception (satellite stem cell count at day 3) is shown in red.}
  \label{fig:calibration}
\end{figure}

\subsection{Perturbations reproduce literature findings}

Across the eight perturbation scenarios, the model qualitatively reproduced published experimental observations. VEGF-A injections produced dose-dependent improvements in CSA recovery and capillary count, with the high (relative concentration $10^3$) injection producing the largest CSA improvement at 28 days and the extra-high ($2 \times 10^3$) injection producing the second largest. Medium and low VEGF-A doses produced transient benefits at day 15 that were not significant relative to control by day 28. Hindered angiogenesis produced a persistent CSA deficit together with elevated ECM collagen density and prolonged macrophage presence at late timepoints, consistent with progression toward fibrosis. IL-10 knockout simulations showed a rise in TNF-$\alpha$ that peaked at day 7 and a compensatory decrease in HGF and TGF-$\beta$, consistent with the expected feedback structure. MCP-1 knockout produced broad increases in most cytokines at 12 hours post-injury, except VEGF-A, followed by a mid-regeneration decrease in TNF-$\alpha$, TGF-$\beta$, IL-10, and MMP-9. Two cases diverged from the experimental literature: TNF-$\alpha$ knockout predicted increased CSA recovery, whereas experiments report decreased recovery, a discrepancy likely attributable to the omission of interferon cross-regulation; and macrophage depletion predicted monotonically decreased TGF-$\beta$ throughout, whereas experimental data show an initial decrease followed by elevation at days 7 and 14, a pattern likely influenced by clodronate liposome pharmacokinetics that are not represented in the model.

\subsection{Nonlinear cytokine-cell relationships}

Beyond reproducing existing experiments, the model generated several nonlinear predictions. VEGF-A dose response was saturating: beyond the high dose, additional VEGF-A did not translate into additional CSA gains, consistent with angiogenic signaling being limited by capillary density or downstream cellular capacity. The relationship between TNF-$\alpha$ and fibroblast count was non-monotonic, with an intermediate concentration supporting maximal fibroblast numbers and higher or lower concentrations suppressing them. HGF concentration was positively correlated with CSA recovery at 28 days up to a threshold, beyond which additional HGF did not improve outcomes. MCP-1 displayed an optimal concentration for M1 macrophage count at day 1, and IL-10 concentration was positively correlated with M2 count at day 3. These patterns illustrate how tissue-scale outcomes arise from multiple overlapping thresholds that are difficult to identify without a spatial, multicellular simulation.

\subsection{Sensitivity analysis and the combined perturbation}

The LHS/PRCC analysis revealed that HGF decay is the most globally influential cytokine parameter: it significantly impacted every model output except M1 cell count. TGF-$\beta$ and MMP-9 decay were positively correlated with CSA recovery, while HGF decay was negatively correlated with CSA, with the correlations strengthening after day 12. Guided by these sensitivities, we designed single-parameter perturbations lowering HGF and VEGF-A decay, raising TGF-$\beta$, MMP-9, and MCP-1 decay, and raising MCP-1 diffusion. All of these except increased MCP-1 decay improved CSA recovery when applied individually. We then simulated a combined perturbation that applied all beneficial changes simultaneously (dropping the counter-productive MCP-1 decay increase). The combined perturbation yielded the highest CSA recovery across all simulations, increased M1 macrophage counts, decreased M2 macrophage counts relative to control, increased fibroblast and SSC peaks beyond any single perturbation, and accelerated capillary regeneration. The combined cytokine perturbation predicted a 13\% improvement in CSA recovery compared with unaltered regeneration at 28 days, exceeding any individual cytokine intervention (Figure~\ref{fig:perturbation}).

\begin{figure}[t]
  \centering
  \includegraphics[width=0.70\linewidth]{figures/fig_perturbation.pdf}
  \caption{CSA recovery at 28 days: combined cytokine perturbation relative to unaltered regeneration. The 13\% improvement is taken verbatim from the experimental log; no other perturbation magnitudes are shown numerically.}
  \label{fig:perturbation}
\end{figure}

\begin{figure}[t]
  \centering
  \includegraphics[width=0.95\linewidth]{figures/fig_sensitivity.pdf}
  \caption{Qualitative LHS/PRCC sensitivity matrix. ``+'' and ``-'' cells indicate statistically significant positive and negative correlations between cytokine decay or diffusion parameters and model outputs at the peak-sensitivity timepoint.}
  \label{fig:sensitivity}
\end{figure}

%========================================================================
\section{Discussion}
\label{sec:discussion}

We developed a spatial ABM of murine skeletal muscle regeneration that combines explicit cytokine diffusion with adaptive microvasculature and a rigorous calibration-validation workflow, and we used the calibrated model to identify combinatorial cytokine perturbations that improve simulated CSA recovery. The central quantitative prediction is that a combined change in cytokine decay and diffusion parameters---lower HGF and VEGF-A decay, higher TGF-$\beta$ and MMP-9 decay, and higher MCP-1 diffusion---produces a 13\% increase in CSA recovery at 28 days relative to unaltered regeneration, which is larger than the effect of any single-parameter intervention considered. Mechanistically, the combined perturbation appears to decouple early and late phases of regeneration. In the early phase, reduced HGF decay prolongs SSC activation while increased TGF-$\beta$ decay and MCP-1 diffusion prevent premature fibroblast over-activation. In the later phase, reduced VEGF-A decay sustains angiogenic signaling and increased MMP-9 decay promotes the anti-inflammatory shift and fiber maturation that close the cascade.

These findings sit naturally alongside experimental reports that combination therapies outperform single-cytokine manipulations. Platelet-rich plasma, which contains VEGF-A, TGF-$\beta$, and other factors, has shown variable clinical benefit in muscle injury, an inconsistency that is plausibly explained by uncontrolled combinatorial composition and timing. Synthetic biomaterials coated with IL-4 have been used to bias macrophage polarization toward M2 in injured muscle, and anti-IL-6 antagonism has improved regeneration in prior work. Our model-based prediction aligns with these experimental directions and points toward specific combinatorial interventions: activation of plasmins to release ECM-bound VEGF-A, HGF-coated biomaterials, TGF-$\beta$ antagonists, nuclear-factor-kappa-B-inhibitory peptides to suppress MMP-9 activity, and recombinant MCP-1 hydrogels to modulate diffusion. Each of these technologies exists in some form, but to our knowledge they have not been combined in a temporally coordinated regimen guided by sensitivity analysis.

The study advances earlier ABMs of muscle regeneration \cite{virgilio2015multiscale,martin2017agent} in three ways. First, cytokine diffusion is explicit and coupled to local collagen density through the Graner-Glazier-Hogeweg lattice representation \cite{graner1992simulation,swat2012compucell3d}. This coupling enables the model to capture interventions that change cytokine transport---for example, plasmin activation that releases bound VEGF-A---that non-spatial models cannot represent. Second, microvasculature is represented by adaptive capillary agents that can become non-perfused in response to nearby necrosis and that sprout new capillaries when local VEGF-A and MMP-9 thresholds are met \cite{carmeliet2003angiogenesis}. This mechanistic capillary layer is essential for the dose-dependent VEGF-A response and the fibrotic signature of hindered angiogenesis that the model reproduces. Third, parameter density estimation via CaliPro \cite{joslyn2021calipro} combined with LHS/PRCC \cite{marino2008sensitivity} provides a systematic calibration-then-exploration workflow that can be applied to other regenerative systems.

The nonlinear patterns identified by the model have direct implications for experimental design. The saturating VEGF-A dose response suggests that studies evaluating angiogenic biomaterials should include doses spanning the threshold identified here; otherwise, dose-ranging experiments risk false negatives. The non-monotonic TNF-$\alpha$--fibroblast relationship suggests that cytokine antagonism strategies should be titrated rather than maximized, and the threshold behavior in HGF--CSA correlations implies that HGF supplementation beyond a critical level may not yield additional benefit, which matters for the design of biomaterial-based HGF delivery systems. More broadly, the model predicts that timing matters: early-phase and late-phase interventions act through different mechanisms, suggesting that a single static cocktail is unlikely to be optimal and that scheduled, staged interventions may outperform bolus delivery. These \emph{in silico} hypotheses can be prioritized against experimental cost, and the model can be used to estimate the effect sizes needed for adequately powered confirmatory experiments.

Several limitations should temper interpretation. First, the model omits several cell types (notably regulatory T cells, eosinophils, and pericytes) and restricts cytokine representation to seven factors. The redundancy evident in the calibration success suggests that these additional factors may have partially overlapping roles, but specific interventions targeting omitted cells cannot be evaluated in the current model. Second, hypertrophy is not represented, so simulated CSA cannot exceed the pre-injury baseline; this restricts applicability to interventions that produce net anabolic remodeling. Third, the simulation is two-dimensional; prior work indicates that 2D and 3D diffusion predictions agree closely for this geometry \cite{hillen2009user}, but spatial heterogeneity in the longitudinal direction cannot be captured. Fourth, calibration draws on male-only experimental data; known sex differences in muscle regeneration \cite{kararigas2022sex} indicate that the model's predictions should not be naively extrapolated to female muscle without retuning. Fifth, parameter density estimation selects one plausible region of parameter space, and other regions may yield qualitatively similar fits with different predictive behavior; exploring parameter ensemble predictions is a natural next step.

Looking forward, we plan to extend parameter density estimation to optimize the timing, composition, and dosing of exogenous cytokine deliveries, using the model as a preclinical design tool. We will also explore injury types and anatomical sites whose microvascular and neural contexts differ from the myotoxin injury used here, and incorporate sex-specific rules as comparable female datasets become available. More broadly, this study illustrates a pattern that may generalize to other regenerative systems: combining a physically grounded spatial framework, systematic parameter density estimation, and global sensitivity analysis turns a descriptive cellular rulebook into a generative source of hypotheses. For muscle injury---where the combinatorial space of possible cytokine interventions dwarfs experimental throughput---this combination offers a principled route from mechanistic understanding to testable therapeutic strategies.

%========================================================================
\bibliographystyle{plain}
\bibliography{references}

\end{document}
